\documentclass{article}
\usepackage[utf8]{inputenc}

\usepackage{amsmath, amssymb, amsthm, amsfonts}
\usepackage{subcaption, enumerate}
\usepackage{bbold}
\usepackage{polynom}
\usepackage{xfrac}
\usepackage{graphicx}
\usepackage{mathtools}
\usepackage{enumitem}

\usepackage{tikz}

\newtheorem{theorem}{Theorem}
\newtheorem{corot}{Corollary}[theorem]
\newtheorem{lemma}{Lemma}
\newtheorem{corol}{Corollary}[lemma]
\newtheorem{proposition}{Proposition}
\newtheorem{corop}{Corollary}[proposition]

\theoremstyle{definition}
\newtheorem{definition}{Definition}
\newtheorem{example}{Example}[definition]
\newtheorem*{notation}{Notation}

\newcommand{\dab}
{
	\begin{tikzpicture}[scale=0.04]
		% head
		\draw[thick] (0, 0) circle (3cm);

		% bent arm
		\draw[thick, rotate around={15 : (0, 0)}] (-3, 0) -- (-3.7, -2) -- (4.1, -1.1);

		% extended arm
		\draw[thick, rotate=20] (3, 0) -- (7, 2);
	\end{tikzpicture}
}
\newcommand{\st}{\text{ s.t. }}

\makeatletter
\renewcommand*\env@matrix[1][*\c@MaxMatrixCols c]{%
  \hskip -\arraycolsep
  \let\@ifnextchar\new@ifnextchar
  \array{#1}}
\makeatother

\usepackage{fullpage}

\title{Determinants Day 4 \\ with Dr Nevard}
\author{Laura Gallo}
\date{October 15, 2020}

\begin{document}
\maketitle

\section{Determinant of Transpose Matrices}
\begin{proposition}
	$\det A=\det A^T$
	\begin{proof}
		We want to prove $\sum_{\sigma\in S_n}(-1)^\sigma\prod_{i=1}^na_{i\sigma_i}=\sum_{\sigma\in S_n}(-1)^\sigma\prod_{i=1}^na_{\sigma_ii}$
		$$\prod_{i=1}^na_{i\sigma_i}=\prod_{\sigma_i=1}^na_{(\sigma^{-1}\sigma_i)\sigma_i}=\prod_{j=1}^na_{\sigma^{-1}_jj}=\prod_{j=1}^na_{\sigma_jj}$$
		$$\therefore\sum_{\sigma\in S_n}(-1)^\sigma\prod_{i=1}^na_{i\sigma_i}=\sum_{\sigma\in S_n}(-1)^\sigma\prod_{i=1}^na_{\sigma_ii}$$
		\dab
	\end{proof}
\end{proposition}

\section{Permanent}
\begin{definition}
	The permanent of a matrix is defined using the formula below:
	$$P(A)=\sum_{\sigma\in S_n}\prod_{i=1}^na_{i\sigma_i}$$
	Interestingly, though the determinant of a matrix is calculated in $\mathcal{O}(n^3)$ time, the permanent of a matrix is calculated in exponential time.
\end{definition}

\section{Recursive Formula for the Determinant}
\begin{equation*}
	\det A =\det \left(
		\begin{matrix}
			a_{11} & a_{12} & \dots & a_{1n} \\
			& & \text{stuff}
		\end{matrix}
	\right) = \sum_{i=1}^n a_{1i} \det \left(
		\begin{matrix}
			0 & \dots & 0 & 1 & 0 & \dots & 0 \\
			& & & \text{stuff}
		\end{matrix}
	\right)
\end{equation*}
\begin{equation*}
	= \sum_{i=1}^n a_{1i} \det \left(
		\begin{matrix}
			0 & \dots & 0 & 1 & 0 & \dots & 0 \\
			& \text{stuff} & & 0 & & \text{stuff}
		\end{matrix}
	\right) = \sum_{i=1}^n a_{1i} (-1)^{i-1} \det \left(
		\begin{matrix}
			1 & 0 & \dots & 0 \\
			0 & & \text{stuff}
		\end{matrix}
	\right)
\end{equation*}
$$=\sum_{i=1}^n a_{1i}(-1)^{i-1}\det A_{1i}$$
where $A_{1i}$ is the $n-1$ by $n-1$ submatrix with the first row and the $i$th column removed. \\
A more general formula is: \\
$$\det A=\sum_{i=1}^n(-1)^{i+j}a_{ij}\det A_{ij}=\sum_{j=1}^n(-1)^{i+j}a_{ij}\det A_{ij}$$

\subsection{Example: Determinant of a 3x3 Matrix}
\begin{equation*}
	\left|
		\begin{matrix}
			2 & 6 & 3 \\
			1 & 9 & 4 \\
			8 & 5 & 5
		\end{matrix}
	\right| =
	2\left|
		\begin{matrix}
			9 & 4 \\
			5 & 5
		\end{matrix}
	\right| - 6\left|
		\begin{matrix}
			1 & 4 \\
			8 & 5
		\end{matrix}
	\right| + 3\left|
		\begin{matrix}
			1 & 9 \\
			8 & 5
		\end{matrix}
	\right|
	= 2(9\cdot5-4\cdot5)-6(1\cdot5-4\cdot8)+3(1\cdot5-9\cdot8)=11
\end{equation*}

\section{Inverses of Matrices}
We have, by multilinearity:
\begin{equation*}
	\sum_{j=1}^n(-1)^{i+j}a_{kj}\det A_{ij}=
	\begin{cases}
		0 & \text{if $i\neq k$} \\
		\det A & \text{if $i=k$} \\
	\end{cases}
\end{equation*}
\\
We can extend this further and say the following: \\
Let $B=(b_{ij})$, where $b_{ij}=(-1)^{i+j}\det A_{ij}$. \\
The $k$th row, $\ell$th column of $AB$ is:
\begin{equation*}
	(AB)_{k\ell}=\sum_{m=1}^na_{km}b_{m\ell}=\sum_{m=1}^na_{km}(-1)^{m+\ell}\det A_{\ell m}
	=\begin{cases}
		0 & \text{if $k\neq \ell$} \\
		\det A & \text{if $k=\ell$} \\
	\end{cases}
\end{equation*}
In other words, $AB=\det A \cdot I$, or $A^{-1}=\frac{B}{\det A}$.
\end{document}

